\section{Núcleo estadístico en Python}
    \label{anexo_python}
    Se reproduce el bloque de \textit{hypothesis\_testing.py} que implementa las
    cinco pruebas con \textit{SciPy} y \textit{statsmodels}. Se omiten la generación
    del conjunto de datos y el código de las figuras 1 a 4 por razones de extensión;
    el script completo está disponible en el repositorio del proyecto.

    \vspace{10pt}

    \lstinputlisting[style=python]{utils/codes/hypothesis_testing_core.py}

    \vspace{10pt}

    La figura interactiva de la Figura \ref{fig_plotly} se construye con
    \textit{Plotly Express}, cuyo argumento \texttt{trendline} delega el ajuste de
    la recta por grupo en \textit{statsmodels}, de modo que las pendientes de la
    Tabla \ref{atenuacion} y las dibujadas en la figura proceden del mismo cálculo.

    \vspace{10pt}

    \lstinputlisting[style=python]{utils/codes/plotly_figure.py}
